%% -*- coding:utf-8 -*-
\chapter{Vorwort}

Dieses Buch verfolgt zwei Ziele: erstens die vergleichende Analyse der wichtigsten syntaktischen Eigenschaften der
\ili{germanisch}en Sprachen und zweitens die Einführung eines bestimmten Formats für die Beschreibung und
den Vergleich von Sprachen. Der Rahmen, in dem die Analysen formuliert werden, heißt \emph{HPSG
  light}. Er basiert auf der Head-Driven Phrase Structure Grammar (HPSG) \citep{ps,ps2,HPSGHandbook} in der speziellen Version, die ausführlich in
\citew{MuellerLehrbuch3} beschrieben wird. HPSG light enthält allerdings keine komplizierten Merkmal"=Wert"=Matrizen
(AVMs). Sofern überhaupt AVMs verwendet werden, sind sie auf das Minimum reduziert und enthalten nur eine kleine Menge
von Merkmalen wie \argst für die Argumentstruktur, \comps für die Komplemente und \spr für den Spezifikator. Alle
übrigen Aspekte der Analysen werden in syntaktischen Bäumen dargestellt, die leichter zu lesen sind. Die Idee
hinter der Einführung von HPSG light besteht darin, ein Werkzeug für Linguistinnen und Linguisten bereitzustellen, die ein
Phänomen detaillierter beschreiben wollen, ohne sich dabei zwangsläufig mit allen
technischen Einzelheiten befassen zu müssen. Der Grad der Formalisierung entspricht dem, was in der Government"=and"=Binding"=Theorie,
im Minimalismus und in den weniger formalen Varianten der Konstruktionsgrammatik üblich ist. Was die eine formale Version
der Konstruktionsgrammatik betrifft, die eine Variante der HPSG ist, nämlich die Sign-Based Construction Grammar (SBCG, \citealp{Sag2012a}), so kann HPSG
light auch als eine Light"=Version von SBCG betrachtet werden, da die Unterschiede in den
abgekürzten Darstellungen und Bäumen, die in diesem Buch verwendet werden, vernachlässigt werden. Die hier vorgestellte Arbeit unterscheidet sich
von nicht"=formaler Arbeit in GB/Minimalismus und Konstruktionsgrammatik in einem wichtigen Punkt: Sie wird
durch implementierte Grammatiken gestützt, die die vollständige Version der HPSG einschließlich einer semantischen Analyse im
Rahmen der Minimal Recursion Semantics (MRS, \citew*{CFPS2005a}) verwenden. Die detaillierten Analysen werden in
Konferenzbänden, Zeitschriftenartikeln und Büchern beschrieben, und die Leserin bzw.\ der Leser ist eingeladen, diese
Quellen zu konsultieren, falls Interesse an den Details besteht. Die implementierten Grammatiken werden
mit der virtuellen Maschine Grammix \citep{MuellerGrammix} verteilt und können von der Webseite des Autors
heruntergeladen werden.\footnote{
\url{https://hpsg.hu-berlin.de/Software/Grammix/}, \today.} Grammix enthält die Grammatiken für
\ili{Deutsch}\footnote{
\url{https://hpsg.hu-berlin.de/Fragments/Berligram/}, \today. Die \ili{deutsch}e Grammatik ist in
\citew{MuellerLehrbuch,MuellerGS} und \citew{MOe2011a,MOe2013a} dokumentiert.
},
\ili{Dänisch}\footnote{
\url{https://hpsg.hu-berlin.de/Fragments/Danish/}, \today. Die \ili{dänisch}e Grammatik ist in
\citew{MOeDanish,MOe2013a,MOe2011a,MOe2013b} dokumentiert.
}, \ili{Englisch}\footnote{
\url{https://hpsg.hu-berlin.de/Fragments/English/}, \today. Die \ili{englisch}e Grammatik ist kleiner als die
\ili{deutsch}e und die \ili{dänisch}e Grammatik. Sie ist ein Proof of Concept einer lexikalistischen Analyse von Passiv, Benefaktiv\-
konstruktionen und Resultativkonstruktionen. Siehe \citew{MuellerLFGphrasal} und \citew{MOe2013a} für Details.
} und \ili{Jiddisch}\footnote{
\url{https://hpsg.hu-berlin.de/Fragments/Yiddish/}, \today. Die \ili{jiddisch}e Grammatik basiert auf
\citew{MOe2011a} und unveröffentlichter Arbeit von Jong-Bok Kim, Alain Kihm und mir zur
Prädikatstopikalisierung im \ili{Koreanisch}en und \ili{Jiddisch}en \citep*{MKK2019a}.
} die im Projekt CoreGram \citep{MuellerCoreGram} entwickelt wurden. Die jeweiligen Webseiten der Grammatiken enthalten eine Liste von Testsätzen,
die von den Grammatiken akzeptiert oder zurückgewiesen werden. Die Leserinnen und Leser sind eingeladen, diese
Sätze in das TRALE-System \citep{DKMM2004a-u,Penn2004a-u} einzugeben, das mit Grammix geliefert wird, und die vollständigen AVMs zu inspizieren.

Das Buch beginnt mit zwei einführenden Kapiteln: Das erste Kapitel stellt die \ili{germanisch}en Sprachen vor
und liefert grundlegende Fakten wie die Zahl der Sprecher, die Gebiete, in denen sie gesprochen werden, und einige historische
Fakten. Das zweite Kapitel behandelt die Phänomene, die im Rest des Buches behandelt werden, \eg Scrambling,
die Stellung von Adverbialen, Passiv, Satztypen und Fernabhängigkeiten. Das dritte Kapitel ist eine
Einführung in Phrasenstrukturgrammatiken, die seit \citegen{Chomsky57a} Formalisierung strukturalistischer Ideen
\citep{Bloomfield33a-u} die Grundlage fast aller Theorien bilden. Kapitel~\ref{chap-psg-xbar} führt nicht nur Phrasenstrukturgrammatiken ein,
sondern auch Grammatiken, die Abstraktionen über Phrasenstrukturregeln verwenden und letztlich zu sehr
abstrakten Grammatiken des auch aus der \xbart
\citep{Jackendoff77a} bekannten Typs führen. Kapitel~\ref{chap-valence} erklärt, wie das Konzept der Valenz
mit abstrakten Phrasenstrukturregeln kombiniert wird, um sicherzustellen, dass die richtige Anzahl und die richtige Art von
Elementen mit einem bestimmten Wort kombiniert wird. Zum Beispiel braucht ein Wort wie \emph{lachen} ein Subjekt und
ein Wort wie \emph{lesen} ein Subjekt und ein Objekt. Das muss irgendwo in einer
Grammatik repräsentiert werden, und Kapitel~\ref{chap-valence} erklärt, wie das in HPSG (light) gemacht wird. Die grundlegenden Unterschiede
in den Analysen von SVO- und SOV-Sprachen werden erläutert. Dieses Kapitel
erklärt außerdem, wie sich die verschiedenen Abfolgen von Subjekt und Objekten in einer Sprache wie dem \ili{Deutsch}en erklären lassen
(das sogenannte \emph{Scrambling}) und wie man die verschiedenen Stellungsmöglichkeiten in Sprachen wie dem \ili{Englisch}en
und den nord\ili{germanisch}en Sprachen einerseits und dem \ili{Deutsch}en, \ili{Niederländisch}en und \ili{Afrikaans} andererseits erfassen kann. Verbalkomplexe werden in Kapitel~\ref{chap-verbal-complex} behandelt, die Verberststellung (zur Fragebildung verwendet) und die Verbzweitstellung (für Aussagen) werden
in Kapitel~\ref{chap-verb-position} erläutert. Passiv und Kasuszuweisung im Allgemeinen werden in
Kapitel~\ref{chap-case} behandelt. Die \ili{germanisch}en Sprachen sind hier besonders interessant, da das \ili{Isländisch}e
zu dieser Sprachgruppe gehört und für seine quirky subjects bekannt ist (Subjekte im Genitiv, Dativ oder
Akkusativ, \citealp{ZMT85a}). Kapitel~\ref{chap-expletives} behandelt Expletivpronomen und wie sie
in den \ili{germanisch}en Sprachen verwendet werden, um Satztypen zu markieren. Zum Beispiel werden Expletiva im
\ili{Deutsch}en in Hauptsätzen verwendet, um die Anfangsposition zu füllen, sodass der Satz eine Aussage ist. Das \ili{Dänisch}e verwendet
Expletiva in eingebetteten Sätzen mit Subjekten als Interrogativelementen. Auch hier beeinflussen die Unterschiede in den
allgemeinen grammatischen Eigenschaften die Grammatik in anderen Bereichen, etwa bei der Stellung von Expletiva.

% The final chapter, Chapter~\ref{chap-HPSG-light}, is for advanced readers. It tries to relate the simplified version
% of HPSG that was used throughout the earlier chapters to HPSG as it is used in full-fledged HPSG publications.

Das letzte Kapitel, Kapitel~\ref{chap-outlook}, fasst kurz zusammen, was im Buch geschehen ist, und
verweist die interessierte Leserin bzw.\ den interessierten Leser auf einige weiterführende Literatur zur HPSG.

\ili{Deutsch}sprachige Folien, die für den von mir mit diesem Buch gehaltenen Kurs entwickelt wurden, sind auf GitHub verfügbar.\footnote{
\url{https://github.com/stefan11/germanic-syntax-slides}, 2021-09-14.
}
\ili{Deutsch}sprachige Vorlesungen zu den Kapiteln finden sich außerdem auf YouTube.\footnote{
\url{https://www.youtube.com/playlist?list=PLXwGGsuPxWRp4AB2LsWH6LKc0II7uc6tg}
}

\section*{Zur Art der Veröffentlichung dieses Buches}

Lehrkräfte an Schulen und an vielen Universitäten werden vom Staat bezahlt, das heißt von der Öffentlichkeit (von Ihnen). Zu
ihren Aufgaben gehört die Erstellung von Lehrmaterial. Es gibt überhaupt keinen Grund, das
Lehrmaterial profit"=orientierten Verlagen zu überlassen. Im Gegenteil, Lehrmaterial sollte offen
und an die Bedürfnisse der Lehrenden anpassbar sein, die es verwenden wollen.

Eine Studie des American Enterprise Institute zeigt, dass der Preis von Hochschulbüchern von 1978 bis 2012 um 812\,\%
gestiegen ist, während die allgemeinen Verbraucherpreise nur um 250\,\% gestiegen sind.\footnote{
\url{https://www.aei.org/carpe-diem/the-college-textbook-bubble-and-how-the-open-educational-resources-movement-is-going-up-against-the-textbook-cartel/}.
2022-12-22.%
} Ähnliche Zahlen gibt es für wissenschaftliche Bücher im Allgemeinen und für Universitätslehrbücher. Mein Lieblingsbeispiel ist ein schmales Lehrbuch
zur Logik, \emph{Logik für Linguisten}, eine Übersetzung des \ili{englisch}en Lehrbuchs \emph{Logic for
Linguists} \citep{AAD73a}. Dieses Buch hat 112 Seiten. Es wurde vom Max Niemeyer
Verlag als Taschenbuch für 9,40€ verkauft. Dieser Verlag wurde von De Gruyter übernommen, und das Buch wird nun für \$126.00/89,95€ als
eBook und \$133,00/94,95€ als gebundene Ausgabe verkauft\footnote{
%  \url{http://www.degruyter.com/isbn/978-3-11-096350-2}. 2014-09-01.
Ich habe 2022 bemerkt, dass De Gruyter dieses Buch nicht mehr anbietet. Ich finde das sogar noch schlimmer.
} (siehe \citealp{MuellerOA} für weitere Beispiele und eine allgemeine Diskussion). Sowohl das eBook als auch das gedruckte Buch sind für Studierende unerschwinglich. Der Ausweg aus dieser höchst
problematischen Situation besteht darin, Bücher im Open Access zu veröffentlichen. Die PDF-Version dieses Buches ist für
alle kostenlos, und die gedruckte Ausgabe ist zu einem angemessenen Preis erhältlich, da das Buch unter
einer Creative-Commons-Lizenz steht und somit nicht einem
profit"=orientierten Verlag gehört und jede und jeder den eigenen Print-on-Demand-Dienst wählen kann, falls
der von Language Science Press standardmäßig angebotene Dienst teurer ist.

\if0
\section*{Gender issues}

\itdopt{M: drop section}

\citet{MB97a} did a study on ten textbooks in syntax published between 1969 and 1994. They showed
that some of the textbooks used examples describing violence against women. Examples like \emph{John
  beats Mary.} or \emph{John hits Mary.} were frequent in the 1970ies and 1980ies. I know of a paper
written by two female authors and one male author containing a \emph{John hits Mary} example. \citet[\page
812]{MB97a} discus some even more extreme cases from the textbooks they examined. Furthermore, females were
depicted in stereotypical situations like teaching children, never in work situations with the
exception of work as secretaries. Some examples were explicitly sexist, referring to women as stupid
and to men as intelligent. 


\citeauthor{MB97a} examined the semantic roles in which women and men appeared, checked for
usage of full NPs, proper names and pronouns. They found that men appear more often than women
in example sentences and they appear more often as agents then as experiencers/patients. While
openly sexist examples disappeared from textbooks in the time since the 80ies, \citet{PCKSDMC2017a}
repeated the studies of \citeauthor{MB97a} and noticed that not much has changed in respect to the
women/men distribution in examples. 

I think that such biased examples contribute to the stereotypes regarding genders that most of us
have since they are deeply entrenched in our societies. Mikaela Wapman and Deborah Belle did a study
asking students of psychology about the following situation: a father and son are in a horrible car
crash and the father gets killed. The son is brought to the hospital and has to undergo an emergency
surgery but the surgeon refuses to do it and say, ``I can’t operate—that boy is my son!''. Question:
How is this possible? People came up with all sorts of explanations like gay fathers, ghost surgeons
and so on, overlooking the possibility of the surgeon being the mother of the child. This shows that
the idea that surgeons are male is deeply entrenched in our societies (the same experiment works in
\ili{German} and it works in the other direction with nurses). The obvious solution to such problems is to
change the employment of women and to change the payment of jobs usually taken by women. This is a
political problem and it would be helpful if women were in at least 50\,\% of the positions in
political systems. Many of the problems that women have can be traced back to a lack of economic
independence. There is a good portray of East \ili{German} women showing the consequences of economic
independence: \url{https://www.mdr.de/zeitreise/schwerpunkte/ostfrauen-106.html}. When the wall came
down and East \ili{German} women and West \ili{German} women united there was much confusion since the issues
they wanted to address were totally different (see also the MDR video documentation). Almost all women in
East Germany worked, child care was available everywhere and so on.

What we seem to have here is an hen--egg problem: society and stereotypes would be different if women
would be in other jobs and better payed, but stereotypes prevent women from even trying because they
feel not welcome in certain jobs. Here is where language plays an important role. \ili{German} marks
gender for professions: one can talk about \emph{Kindergärtner} `male nurse' and
\emph{Kindergärtnerinnen} `female nurse'. Looking at the \ili{German} grammatical system, the plural is unmarked for
gender, but there is the form \emph{Kindergärtnerinnen} that makes gender explicit. Studies show
that recipients consider women less often, if the unmarked form is used. For instance, \citet{SSB2001a}
showed that people come up with more names of male musicians, when asked for their favorite \emph{Musiker} rather than
their favorite \emph{Musiker} `musician' or \emph{Musikerin} `female musician'. One proposal solved this
problem by using abbreviations for fusing both variants: \emph{MusikerIn} `male or female musician'. Further more inclusive
forms have been developed including people with diverse genders: \emph{Musiker*in}. These forms are
standard in the communication of the \ili{German} Research Foundation and some universities (\eg the
Humboldt-Universität zu Berlin, where I am currently working).    

While the so-called \emph{Binnen-I} helps to make women visible in communication this does not help
as far as textbooks are concerned since respective forms are rather rarely used in examples. So, to
make non-male readers feel more welcome, other means have to be employed. One has to work on the
referents that are used in examples. The names used most frequently in examples in the papers/textbooks from which I was learning syntax
(written in the 80ies and 90ies) were \emph{John} and \emph{Mary} or \emph{Karl} and \emph{Maria} in
\ili{German} publications. (Books on HPSG using gender neutral names like \emph{Kim} and \emph{Sandy} were an early exception \citep{ps,ps2}.)
I have never used discriminating examples in my publications and I avoided stereotypes, but I also used
\emph{Karl} and \emph{Maria} and I had \emph{den Mann} `the man' reading books and \emph{der Mann}
`the man' giving books to \emph{der Frau} `the woman' or \emph{dem Kind} (neuter, `the child'). One
reason for having more male referents in \ili{German} examples is that the masculine inflection paradigm
differentiates between all four cases while the feminine one collapses nominative and accusative as
well as genitive and dative. The neuter is somewhere in the middle: nominative and accusative are syncretic. Table~\ref{table-German-case-syncretism} gives an overview of the situation.
\begin{table}
\begin{tabular}{llll}\lsptoprule
           & the man    & the woman & the child\\\midrule
nominative & der Mann   & die Frau  & das Kind\\
genitive   & des Mannes & der Frau  & des Kindes\\
dative     & dem Mann   & der Frau  & dem Kind\\
accusative & den Mann   & die Frau  & das Kind\\
\lspbottomrule
\end{tabular}
\caption{\label{table-German-case-syncretism}Inflectional marking within noun phrases of feminine and masculine gender. Only the masculine
  paradigm is unambiguously marked for case.}
\end{table}  
If one wants to avoid male referents, one could use animals, but depending on the verb one
uses, a dolphin as subject would not make sense. So, one way to solve this problem, is to use proper
names like \emph{Kim}, \emph{Sandy}, and \emph{Chris} that are unisex. This is of course only
possible if case is unimportant since proper names do not inflect for case in almost all of the \ili{Germanic}
languages (\ili{Icelandic} being the exception, \citealp[\page 443]{ZMT85a}). 

As mentioned above, there is some tradition of what names are used in linguistic examples: for
\ili{English} it was \emph{John} and \emph{Mary}, for \ili{German} \emph{Karl} and \emph{Maria}, \ili{Dutch} authors
used \emph{Jan} and \emph{Piet} and \emph{Marie}. Since this book is about \ili{Germanic} languages I
thought I collect names that are typical for the languages under discussion. It is clear that the
names trick does not work globally. I learned that \emph{Gert} can be used both for women and
men in the \ili{Scandinavian} countries, but in Germany this name exclusively refers to
males.\footnote{%
See Bernadette La Hengst's song \emph{Ein Mädchen namens Gerd} `A girl named Gerd', which is a parallel to \emph{A boy
  named Sue} by Johnny Cash: \url{https://www.youtube.com/watch?v=Y6WVcOk4D3c}, 2020-06-26.
} In
addition speaker communities may associate certain names with women and girls since there are more females
than males with this name in a certain community and vice versa. And finally there are differences
on the individual level: you may happen to know three female Connys and no male one but this may
change during the course of your live. So an example about Conny reading a book or feeding a child
could work for or against a certain stereotype, depending on the experiences of the reader. And of course the assignment to a certain gender may change in
time: a name that is given more often to girls may be given more often to boys some decades later. 
So, what I did is not perfect but since I state here that all names are meant to be unisex including
non-binary people, it is a step into the right direction.

In a discussion on twitter, XY pointed out to me that names do not necessarily fit the gender of the
referent. There are famous examples of a male US American lawyer with the name Sue\footnote{
\url{https://en.wikipedia.org/wiki/Sue_K._Hicks}, 2020-07-06.
%\emph{Johnny Cash Is Indebted to a Judge Named Sue}. The New York Times, July 12, 1970, p.\ 66.
} and of women with the name Michael (the columnist Michael Sneed, Chicago Sun-Times, and The
Bangles' bassist Michael Steele).\footnote{
  \url{https://en.wikipedia.org/wiki/Michael_Burnham}, 2020-07-06.
}
So in principle
one would have to make the gender of the referent explicit. But this is exactly what I want to avoid
by using unisex names: I do not want to entrench stereotypes. By not using \emph{Sue}, I do not have to say
that Sue may be a man.

A special case are non-binary people. They may not feel represented in examples with unisex
names like \emph{Kim} and \emph{Sandy}. This is the reason why they sometimes choose names that are
not used by binary people. Kirby Conrod suggested that non-binary people give permission to use their names and allowed me to
use theirs (Kirby's), which I do at some places in the book. 
% https://nonbinary.miraheze.org/wiki/Names

%Kim Gerdes (\ili{German} linguist), Kim Wilde (\ili{British} pop singer).

Table~\ref{table-given-names-used-in-the-book} shows the names I am using throughout the book.

\begin{table}
\begin{tabular}{llp{4cm}p{4cm}}\lsptoprule
language    & name      & female basis         & male basis\\\midrule
\ili{Danish}      & Gert/Gerd & Asgjerd, Hallgerd, Hildegard, Ingjerd        & Gerhard\\
\ili{Dutch}       & Robin     &                      & Robrecht, Robert\\ %https://nl.wikipedia.org/wiki/Robin_(voornaam)
            & Benedikt \\
\ili{German}      & Conny     & Constanze, Cornelia  & Conrad, Constantin\\
            & Aicke     & & Eckehard \\
\ili{English}     & Chris     & Christina, Christine & Christian, Christoph(er)\\
            & Kim       & Kimberly             & Kimberly, Kimball,\newline Joachim, Joakim\\
            & Sandy     & Sandra               & Alexander, Sander, Alasdair\\
non-binary  & Kirby\\
            & Elvan\\% from Wiki
\lspbottomrule
\end{tabular}
\caption{\label{table-given-names-used-in-the-book}This table lists the names used in the book for
  the languages given}
\end{table}
%\end{sideways}

\fi

~\medskip

\noindent
Berlin, 21.~April 2023\hfill Stefan Müller
%Berlin, \today\hfill Stefan Müller



%% -*- coding:utf-8 -*-
\section*{Vorwort für die zweite Auf\/lage}



%      <!-- Local IspellDict: de_DE -->
